\documentclass[10pt]{article}

\usepackage[utf8]{inputenc}

\usepackage[T1]{fontenc}

\usepackage{amsmath}

\usepackage{amsfonts}

\usepackage{amssymb}

\usepackage[version=4]{mhchem}

\usepackage{stmaryrd}

\usepackage{graphicx}

\usepackage[export]{adjustbox}

\graphicspath{ {./images/} }

\usepackage{bbold}



\begin{document}

с франц., 2 изд., М., 1958; [6] 3 а м а н с к и й М., Введение в современную алгебру и анализ, пер. с франц., М., $1974 ;$ [7] M., 1981 ; [8] Х а у с д о р ф Ф., Теория множеств, пер. с нем. M.- J., 1937 . П. Д. Кудрлвчев.



ПРЕДЕЛЬНАЯ ТОЧКА м н о ж е с т в а - точка, в любой окрестности к-рой содержится по крайней мере одна точка данного множества, отличная от нее самой. Рассматриваемые множества и точка предполагаются принадлежащими нек-рому топологич. пространству. Множество, содержащее все свои П. т., наз. з а м к н у т ы м. Совокупность всех П. т. множества М наз. п рои з во д ны м м но же с тво м и обозначается $M^{\prime}$. Если рассматриваемое топологич. пространство $X$ удовлетворяет первой аксиоме отделимости (для любых двух его точек $x$ и $y$ существ ует окрестность $U(x)$, не содержащая точку $y)$, то каждая окрестность П. т. нек-рого множества $M \subset X$ содержит бесконечно много точек этого множества и производное множество $M^{\prime}$ - замкнуто. Всякая прикосновения точка множества $M$ является либо его П. т., либо изолированной.



Лит.: [1] А л е к с а н д р о в П. С., Введение в теорию множеств и общую топологию, М., 1977; [2] Х а у с д о р ф Ф., Теория множеств, пер. с нем., М.- Л., 1937.



ПРР. Д. Кудрлвчев. ПРЕДЕЛЬНАЯ ТОЧКА т р а е к т о ри и $\left\{f^{t} x\right\}$



\begin{center}

\includegraphics[max width=\textwidth]{2023_07_08_37c8fe0955bf5d32875dg-1(1)}

\end{center}



\[

x_{\alpha}=\lim _{k \rightarrow \infty} f^{t^{\prime}} x

\]



$(\alpha-\Pi$ р е д ль ь а я т о ч к а) или



\[

x_{\omega}=\lim _{k \rightarrow \infty} f^{t^{k} x}

\]



( следовательность такая, что $t_{k} \rightarrow-\infty \quad$ при $k \longrightarrow \infty$ в (1) или $t_{k} \longrightarrow+\infty$ при $k \longrightarrow \infty$ в (2) и пределы (1) или (2) существуют.



Для траектории $\left\{f^{t} x\right\}$ динамич. системы $f^{t}$ (или, иначе, $f .(t, x)$, см. [1]) $\alpha-\Pi$. т. $(\omega-\Pi$. т.) - то же, что $\omega-\Pi$. т.



\includegraphics[max width=\textwidth, center]{2023_07_08_37c8fe0955bf5d32875dg-1}

темы с обращением времени). Множество $\Omega_{x}\left(A_{x}\right)$ всех $\omega-\Pi$. т. ( $\alpha-\Pi$. т.) траектории $\left\{f^{t} x\right\}$ наз. $\omega-\Pi$ р ө д е л ьны м $(\alpha-\Pi$ р е д л в ны м) мно же с т в о м этой траектории



Лum.: [1] Н е мыции й В. В., С т е п а но в В. В., Качественная теория дифференциальных уравнений, 2 изд., М., 1949. ПДЕЛЬНОГО ПОГЛОШЕНИЯ М. ПРЛлионииков. способ однозначного выделения решений уравнений, аналогичных Гельмгольца уравнению, с помощью введения бесконечно малого поглощения, Математич. смысл П. П. П. состоит в следующем. Пусть $\Omega-$ неограниченная область в $\mathbb{R} n, P$ - самосопряженный опөратор в $L_{2}(\Omega)$, задаваемый дифференциалыным выражөнием $P\left(x, \frac{\partial}{\partial x}\right), x \in \Omega$, п однородными граничными условиями на $\partial \Omega, \lambda$ - точка непрерывного спектра оператора $P$. Тогда при $\varepsilon \neq 0$ уравнение



\[

P u_{\varepsilon}=(\lambda+i \varepsilon) u_{\varepsilon}+f

\]



однозначно раз решимо в $L_{2}(\Omega)$ и в нек-рых случаях можно выделить решения $u=u_{ \pm}$уравнения



\[

P u=\lambda u+f

\]



с помощью предельного перехода



\[

u_{ \pm}=\lim _{\varepsilon \rightarrow \pm 0} u_{\varepsilon}

\]



При әтом предполагается, что $f$ имеет компактный носитель, а сходимость $u_{\varepsilon} \rightarrow u_{ \pm}$при $\varepsilon \rightarrow \pm 0$ понимается в смысле $L_{2}\left(\Omega^{\prime}\right)$, где $\Omega^{\prime}$ - произвольная ограниченная область в $\Omega$. Так как $\lambda$ - точка спектра оператора $P$, то указанный предел в $L_{2}(\Omega)$, вообще говоря, нө существ YeT.



Впервые П. П. П. был сформулирован для уравнөния Гельмгольца в $\mathbb{R}^{2}$ (см. [1]):



\[

\begin{gathered}

\left(\Delta+k^{2}\right) u=-f, \Omega=\mathbb{R}^{2}, \\

P=-\Delta, \lambda=-k^{2}<0 .

\end{gathered}

\]



Выделяемые с помощью этого принципа решения $u_{ \pm}$ соответствуют расходящимся или сходящимся волнам и удовлетворяют излучения условию на бесконечности. Эти результаты были перенесены (см. [2] , [3]) на әллиптические краевые задачи во внешности ограниченной области в $\mathbb{R} n$ для оператора



\[

\boldsymbol{P}\left(x, \frac{\partial}{\partial x}\right)=-\sum_{k, j=1}^{n} \frac{\partial}{\partial x_{k}}\left(a_{k j} \frac{\partial}{\partial x_{j}}\right)+q(x),

\]



где коэффициенты $a_{k j}(x)$ достаточно быстро стремятся к константам при $|x| \rightarrow \infty$. Для справедливости П. П. ІІ. в этом случае нөобходимо трөбовать, чтобы $\lambda$ не было собственным значөнием оператора $\boldsymbol{P}$ или $f$ была ортогональна собствендым функциям. Төорема Като (см. [3]) даөт достаточные условия отсутствия собственных значений на непрерывном спектрө оператора $P=-\Delta+q(x)$ Такая теорема получена для оператора (*) (см. [3]). П. П. П. обоснован для нек-рых областей с некомпактной границей (см. [3], [4]).



П. П. п. и соответствующие условия излучения найдены для уравнений любого порядка и систем уравнений (см. [5], [6]); они состоят в следующем.Пусть $P\left(i \frac{\partial}{\partial x}\right)$ эллиптический (или гипоэллиптический) оператор, удовлетворяющий условиям: 1) многочлен $P(\sigma)$ имеет действитөлыные коэффициенты, 2) поверхность $\boldsymbol{P}(\sigma)=0$, $\sigma \in \mathbb{R}^{n}$, распадается на $\varkappa$-связных гладких поверхностей $S_{j}$ с отличной от нуля кривизной, 3) grad $P(\sigma) \neq 0$ на $S_{j}$. Пусть на $S_{j}$ заданы ориентации, т. е. независимо для каждой поверхности выбраны направлөния нормали $v$. Пусть $\omega=\frac{x}{|x|}, \sigma_{j}=\sigma_{j}(\omega)$ - точка на $S_{j}$, в к-рой $v$ и $\omega$ имеют одинаковые направления, и $\mu_{j}(\omega)=\left(\sigma_{j}(\omega)\right.$, $\omega)$. Тогда функция $u(x)$ удовлетворяөт условиям излучения, еслй она представима в виде



\[

\begin{gathered}

u=\sum_{j=1}^{x} u_{j}(x), u_{j}=O\left(r^{(1-n) / 2}\right), \\

\frac{\partial u_{j}}{\partial r}-i \mu_{j}(\boldsymbol{\omega}) u_{j}=0,\left(r^{(1-n) / 2}\right), r \rightarrow \infty .

\end{gathered}

\]



Эти условия выделяют единственное решение уравнения



\[

P\left(i \frac{\partial}{\partial x}\right) u=f, x \in \mathbb{R}^{n},

\]



для любой функции $f$ с компактным носителем. П. П. П. для этого уравнения заключается в том, что әто же решение получается в предөле при $\varepsilon \rightarrow+0$ из однозначного определяемого решения $u_{\varepsilon}(x) \in L_{2}(\mathbb{R} n)$ әллиптич. уравнения



\[

P\left(i \frac{\partial}{\partial x}\right) u_{\varepsilon}+i \varepsilon Q\left(i \frac{\partial}{\partial x}\right) u_{\varepsilon}=f

\]



где $Q(\sigma)$ имеет действительные коәффициенты и $Q(\sigma) \neq$ $\neq 0$ на $S_{j}$. В зависцмости от набора $\operatorname{sign}_{\sigma \in S_{j}} Q(\sigma), 1 \leqslant j \leqslant x$, в пределе получаются решения с условиями излучения, соответствующими той или иной ориентации поверхностей $S_{j}$. Этот принцип обоснован для уравнений и систем любого порядка с переменными коэффициентами во внешности ограниченной области (см. [5], [6]), а также в случае невыпуклых $S_{j}$; для этих уравнений имеется и теорема единственности типа Като.





\end{document}